% HiMCM 2021 Problem A --- Task 5 homeowner letter (~1.5 pages)
% Prose: ~80% ASD-STE100 (short sentences, one idea, active voice).
\documentclass[11pt]{letter}
\usepackage[margin=0.9in]{geometry}
\usepackage{setspace}
\usepackage[hidelinks]{hyperref}
\pagestyle{empty}
\setstretch{1.08}

\begin{document}
\begin{center}
{\Large\bfseries Powering Your Independent Home:}\\[0.2em]
{\large A Practical Guide to Sizing Solar and Battery Systems}\\[0.45em]
{\normalsize HiMCM 2021 Problem A \quad Team XXXX \quad (plain-language brief)}
\end{center}
\vspace{0.6em}

\noindent
Dear homeowner,

\noindent
If you want lights, a fridge, and a Wi-Fi router when the utility is
gone, you run a tiny power company. This letter is the non-math version
of our contest study. The numbers below come from our Phoenix-climate
year and a \$10{,}000 hardware budget. They do not come from a sales
brochure.

\medskip
\noindent\textbf{1. Do not stack more solar panels by habit. Storage is
the lifeline.}
Panels make energy only while the sky is bright. In our year the array
produced about \(12{,}300\,\mathrm{kWh}\). The house wanted about
\(15{,}400\,\mathrm{kWh}\). No pack can invent the missing
\(3{,}100\,\mathrm{kWh}\). Extra panels shrink the annual energy gap.
Packs cover the nights and the cloudy streaks. If you buy modules and
skip usable storage, you still sit in the dark at 8\,p.m.\ in February.
Spend the next dollar on kilowatt-hours you can hold overnight. Do not
spend it on a fifth row of glass that idles after sunset.

\medskip
\noindent\textbf{2. How to read a pack nameplate: usable kWh and
depth of discharge.}
Advertised ``13.5\,kWh'' is not what you may drain every day. Lithium
iron phosphate packs in our catalog keep a safety window (about 10\% to
95\% state of charge). Lead-carbon is stricter (about 30\% to 90\%).
The mix we could buy under \$10{,}000 is three small LFP packs plus one
lead-carbon pack (\(24.6\,\mathrm{kWh}\) nameplate). It behaves as empty
at about 18\% indicated SOC. That is chemistry protection. It is not a
software bug. When a salesperson says ``zero to full every day,'' ask:
\emph{what is the warrantied usable window, round-trip efficiency,
continuous kilowatts, and cycle count to 90\% remaining capacity?}
A 5\,kWh pack you may cycle 2{,}500 times is not the same product as a
9.6\,kWh lead pack with 1{,}200 cycles. Price per \emph{usable}
kilowatt-hour. Do not price the screen.

\medskip
\noindent\textbf{3. What \$8{,}000 vs.\ \$10{,}000 actually bought.}
At \$8{,}000 the model bought two small LFP packs and the lead-carbon
pack. At \$10{,}000 it added a third LFP pack (total capex \$9{,}700).
A name-brand 13.5\,kWh wall pack listed at \$10{,}500 did \emph{not}
fit the \$10{,}000 cap. Reliability improved. The year did not become
zero-outage. Thousands of hours still had some unserved load, because
the annual solar crop is smaller than the annual appetite. Treat
\$8k--\$10k as a first staircase. It is not a miracle.

\medskip
\noindent\textbf{4. Three household rules for a dark, cold stretch.}
Our household meeting has three roles: a parent who wants the freezer
alive, a remote worker who needs the laptop and router, and a butler
who watches the 20\% SOC line. The meeting always returns the same
rules:
\begin{enumerate}
  \item \textbf{Keep the survival circuit.} Fridge, well pump, router,
        work computer. Unplug those ``to save the pack'' and you lose
        the food and the job.
  \item \textbf{Park the wet room.} Dryer and dishwasher are optional
        for 48 hours. Laundry can wait. Insulin cannot.
  \item \textbf{Nudge the thermostat. Do not kill the house.} A couple
        of degrees cooler (or warmer in a heat wave) cuts heat-pump
        draw. You still do not live at outdoor temperature.
\end{enumerate}
In a 72-hour February test, that etiquette trimmed about
\(23\,\mathrm{kWh}\) of deferrable use. Hours with unmet demand fell
from 16 to 5. It did not refill an already-empty pack. Smart shedding
is a shock absorber. It is not a second solar farm.

\medskip
\noindent\textbf{5. Look ahead. Do not gamble the mortgage.}
Sodium-ion packs look cheaper per delivered kWh in a 20-year
spreadsheet (our scenario: about \$0.16/kWh versus about \$0.41/kWh for
today's LFP wall pack). They are not yet a walk-in retail replacement.
Cement that stores a trickle of electricity in a wall is a research
story (very low energy per kilogram). Buy a warrantied LFP or similar
pack you can service in 2026. Keep the rest on a watch-list. Do not
write a purchase order.

\medskip
\noindent
Live within the energy you harvest. Size the pack by \emph{usable}
kilowatt-hours. Agree---before the storm---which plugs are sacred.
That is independence you can explain at the kitchen table.

\vspace{0.9em}
\noindent
Sincerely,\\
Team XXXX \quad Off-grid solar and storage study
\end{document}
